\section{Preguntas y respuestas destacadas: perspectivas clave en las preguntas profundas}

\subsection{Origen del sesgo en los métodos estructurales}

\textbf{P}: ¿Por qué la distribución de generación de los métodos puramente estructurales presenta un sesgo severo?

\textbf{R}: La causa fundamental es el \textbf{sesgo de muestreo de los datos estructurales}. Las proteínas con estructura resuelta experimentalmente (alrededor de 300,000 en la base de datos PDB) tienden a presentar las siguientes características:
\begin{itemize}
    \item \textbf{proteínas globulares} (globular proteins): compactas, solubles en disoluciones acuosas, fáciles de manejar experimentalmente
    \item \textbf{ricas en $\alpha$-helix}: este tipo de elemento estructural está sobrerrepresentado entre las proteínas ya resueltas
\end{itemize}

El artículo de EvoDiff aporta datos detallados al respecto: las proteínas generadas por EvoDiff cubren mejor otros elementos estructurales además de la $\alpha$-helix (como la $\beta$-sheet, entre otros), mientras que los métodos puramente estructurales tienden a sobremuestrear proteínas ricas en $\alpha$-helix.

\subsection{Perspectivas de la fusión de secuencia y estructura}

\textbf{P}: ¿Podrá EvoDiff integrar información estructural en el futuro?

\textbf{R}: Sí, esta es una dirección de investigación activa en el campo del aprendizaje automático aplicado a proteínas. Las estrategias concretas de fusión incluyen:
\begin{itemize}
    \item mapear la secuencia y la estructura a un \textbf{espacio de representación común}
    \item construir un \textbf{modelo generativo conjunto de secuencia y estructura}
    \item usar la información estructural para \textbf{restringir y guiar} la generación de la secuencia
\end{itemize}

Las dos modalidades de datos son \textbf{complementarias}: los datos estructurales son escasos en volumen pero aportan restricciones geométricas precisas, mientras que los datos de secuencia son abundantes en volumen y cubren una amplia diversidad biológica. Fusionarlos promete reunir las ventajas de ambos.

\subsection{Resumen de la sección}

La sesión de preguntas y respuestas reveló dos perspectivas técnicas importantes: (1) el sesgo sistemático de los datos estructurales es la limitación fundamental de los métodos puramente estructurales; (2) la fusión de secuencia y estructura es el camino necesario para lograr capacidades de diseño de proteínas más sólidas.

